\documentclass[10pt]{beamer}

\usetheme[progressbar=frametitle]{metropolis}
\usepackage{appendixnumberbeamer}
\usepackage{smartdiagram}
\usepackage{booktabs}
\usepackage[scale=2]{ccicons}
\usepackage{tabto}
\usepackage{hyperref}
\usepackage{pgfplots}
\usepgfplotslibrary{dateplot}

\usepackage{xspace}
\newcommand{\themename}{\textbf{\textsc{metropolis}}\xspace}

\title{Vershon control}
\subtitle{(using git)}
%\date{\today}
%\date{}
\author{\textbf{Danny Awty-Carroll}}
\institute{\tiny{\date{\today}}}
\titlegraphic{\hfill\includegraphics[height=1cm]{Aberystwyth_University_logo}}


\begin{document}

\maketitle



\begin{frame}{Sections}
This will focus on using git in windows with a UI
\newline

  \setbeamertemplate{section in toc}[sections numbered]
  \tableofcontents[hideallsubsections]
\end{frame}




\section{Why version control?}


\begin{frame}[fragile]{What is version control}
This is just the manigment of vershons of a document.\\
One document through time.
  \begin{columns}[T]
    \begin{column}{.5\textwidth}
	\includegraphics[width=6cm]{Figs/old/no}
    \end{column}
    \begin{column}{.5\textwidth}
	\includegraphics[width=6cm]{Figs/old/date}
    \end{column}
  \end{columns}
All of us use some version control
\end{frame}


\begin{frame}[fragile]{Where things get complicated}
Numbering or dating documents works OK but can fall down when\\
	\begin{itemize}
		\item There are multiple documents that need to work together (i.e. a script and data)
		\item There are multiple people working on the documents (are they on the latist vertion and merging changes)
		\item There are updates to the progect that may brake things
		\item Line changes need to be reviewed
	\end{itemize}

Some of this is solved in programs like word with track changes (so you see who altered what when)
\end{frame}


\begin{frame}[fragile]{What version control systems can do}
	\begin{itemize}
		\item Give line by line user by user line changes
		\item Make a stucher to vershons so there can be side branches
		\item Minimise problems of multiple users working on the same document at the same time
	\end{itemize}
	\includegraphics[width=11cm]{Figs/GHD/change}
\end{frame}


\begin{frame}[fragile]{What version control systems can't do}
	\begin{itemize}
		\item Add much extra control to binary files (not plan text)
		\item Back up in real time
	\end{itemize}
	\includegraphics[width=11cm]{Figs/GHD/binarychange}
\end{frame}



\section{Version control system}


\begin{frame}[fragile]{What version control system?}
The first thing about version control systems is which system (we will cover git) \\
  \begin{columns}[T]
    \begin{column}{.8\textwidth}
	\smartdiagram[bubble diagram]{Effective\\ version\\ control,
 	System, Storage/\\platform, UI/\\interface}
    \end{column}
    \begin{column}{.2\textwidth}
	\includegraphics[width=2cm]{Figs/Git-logo}
    \end{column}
  \end{columns}
\end{frame}


\begin{frame}[fragile]{What version control system?}
  \begin{columns}[T]
    \begin{column}{.8\textwidth}
	\begin{itemize}
		\item There are two populer version control systems (git  SVN)
		\item We will cover git as it is the most populer, the one I know and easyist to use
		\item Git was developed in 2005 by Linus Torvalds (Principal developer of Linux)
  		\includegraphics[width=2cm]{Figs/Linus_Torvalds}
	\end{itemize}
    \end{column}
    \begin{column}{.2\textwidth}
	\includegraphics[width=2cm]{Figs/SVN} \newline \newline
	\includegraphics[width=2cm]{Figs/Git-logo}
    \end{column}
  \end{columns}
\end{frame}


\begin{frame}[fragile]{What is git?}
  \begin{columns}[T]
    \begin{column}{.8\textwidth}
    \end{column}
    \begin{column}{.2\textwidth}
	\includegraphics[width=2cm]{Figs/Git-logo}
    \end{column}
  \end{columns}
	\begin{itemize}
		\item Git is just system of version control
		\item By default it is used through a terminal
		\includegraphics[width=4cm]{Figs/git/status}
		\item Git has progects cauld reposetrys (or repos)
		\item It then has a tree struchure to manige the vershons
	\end{itemize}
\end{frame}


\begin{frame}[fragile]{Tree structure?}
  \begin{columns}[T]
    \begin{column}{.8\textwidth}
	\includegraphics[width=9cm]{Figs/git/tree}
    \end{column}
    \begin{column}{.2\textwidth}
	\includegraphics[width=2cm]{Figs/Git-logo} \newline  \newline
	The master branch is in gray with colour branches coming of and then being merged back
    \end{column}
  \end{columns}
\end{frame}


\begin{frame}[fragile]{Key comards?}
  \begin{columns}[T]
    \begin{column}{.8\textwidth}
	In the terminal interface there are some useful commards:
    \end{column}
    \begin{column}{.2\textwidth}
	\includegraphics[width=2cm]{Figs/Git-logo}
    \end{column}
  \end{columns}
	\begin{itemize}
		\item \textit{init} - inisholies a repo
		\item \textit{status} - tells you if you have flies out of sync with the current vertion
		\item \textit{commit} - adds commits to the current vertion - a change with a comment and an ID
		\item \textit{add} - adds new files and current commit
		\item \textit{remove} - removes files and current commit
		\item \textit{push} - pushes comits to an online repo (only if using a git platform)
		\item \textit{pull} - pulls commits from an online platform (only if using a git platform)
	\end{itemize}
\end{frame}



\section{Platform}


\begin{frame}[fragile]{Which platform?}
You can just use git with the treminal on your PC or you can store work on an online reposetry mainging site, the most populer of these are GitHub and Bitbucket\\
  \begin{columns}[T]
    \begin{column}{.8\textwidth}
	\smartdiagram[bubble diagram]{Effective\\ version\\ control,
 	System, Storage/\\platform, UI/\\interface}
    \end{column}
    \begin{column}{.2\textwidth}
	\includegraphics[width=2cm]{Figs/git/Octocat} \newline  \newline   \newline 
	\includegraphics[width=2cm]{Figs/git/Bitbucket}
    \end{column}
  \end{columns}
\end{frame}


\begin{frame}[fragile]{Why not just store the repostry on dropbox?}
You can and it would be baked up and sherable, but there are problems:\\
  \begin{columns}[T]
    \begin{column}{.8\textwidth}
	\begin{itemize}
		\item You can run into Dropbox syncing and git version control clashes
		\item If you shere the files with a colaborator they are't identifyed separately to you
		\item You or a colaborator can completely mess up the git system (mostly by deleating the records in the `.git' folder)
		\item If you do you can't just re download
	\end{itemize}
    \end{column}
    \begin{column}{.2\textwidth}
	\includegraphics[width=2cm]{Figs/git/Dropbox}
    \end{column}
  \end{columns}
\end{frame}


\begin{frame}[fragile]{Which platform?}
  \begin{columns}[T]
    \begin{column}{.8\textwidth}
	\textbf{Bitbucket}
	\begin{itemize}
		\item Alows unlimetd public or privet repostorys
		\item Charges per colcolaberaters over 5 on the repostorys 
		\item 1GB stroige for large flies all repostorys, 5GB if using a .ac email
	\end{itemize}
	\textbf{GitHub}
	\begin{itemize}
		\item Has unlimeted public repostorys 
		\item With unlimeted collaborators
		\item Pay for privet repostorys
		\item However as a student or academic you can sign up for free privet repostorys
		\item Github is about $\times$4 more populer than Bitbucket
	\end{itemize}
    \end{column}
    \begin{column}{.2\textwidth}
	\includegraphics[width=2cm]{Figs/git/Bitbucket} \newline  \newline   \newline \newline 
	\includegraphics[width=2cm]{Figs/git/Octocat}
    \end{column}
  \end{columns}
\end{frame}



\section{UI}


\begin{frame}[fragile]{Which UI?}
You can just use git with the treminal but there is a frendlyer interfase\\
  \begin{columns}[T]
    \begin{column}{.8\textwidth}
	\smartdiagram[bubble diagram]{Effective\\ version\\ control,
 	System, Storage/\\platform, UI/\\interface}
    \end{column}
    \begin{column}{.2\textwidth}
	\includegraphics[width=2cm]{Figs/git/terminal} \newline  \newline   \newline 
	\includegraphics[width=2cm]{Figs/git/gitdesktop} \newline  \newline   \newline 
	\includegraphics[width=2cm]{Figs/git/Sourcetree}
    \end{column}
  \end{columns}
\end{frame}


\begin{frame}[fragile]{Which UI?}
  \begin{columns}[T]
    \begin{column}{.8\textwidth}
  	\textbf{The Terminal\\}
	You can just use the termainal with the raw git commands as seen before... but:
	\begin{itemize}
		\item If you are not using it enyway or are happy using it a graphical UI is nice 
		\item You don't need to remember commands or reposty names 
	\end{itemize}
	\includegraphics[width=6cm]{Figs/git/terminaluse}
    \end{column}
    \begin{column}{.2\textwidth}
	\includegraphics[width=2cm]{Figs/git/terminal} \newline  \newline  \newline 
    \end{column}
  \end{columns}
\end{frame}



\begin{frame}[fragile]{Which UI?}
  \begin{columns}[T]
    \begin{column}{.8\textwidth}
  	\textbf{Git Desktop\\}
	The easyist to use and most populer UI is git desktop:
	\begin{itemize}
		\item This is made by github and intigrates very well with github or local (on pc) reposotrys
		\item It will work with other hosting platforms but less easly
		\item Is the easyist to use
		\item Will install git for widows \url{https://git-scm.com/download/win}
	\end{itemize}
	\includegraphics[width=8cm]{Figs/GHD/outline_02}
    \end{column}
    \begin{column}{.2\textwidth}
	\includegraphics[width=2cm]{Figs/git/gitdesktop} \newline  \newline  \newline 
    \end{column}
  \end{columns}
\end{frame}


\begin{frame}[fragile]{Which UI?}
  \begin{columns}[T]
    \begin{column}{.8\textwidth}
  	\textbf{Sourcetree\\}
	The hosting platform indipendet:
	\begin{itemize}
		\item Sourcetree is uned by atlassian who ouns Bitbucket
		\item It will work with other hosting platforms easly
		\item It has more complex controlls
		\item Will install git for widows \url{https://git-scm.com/download/win}
	\end{itemize}
	\includegraphics[width=8cm]{Figs/ST/ST_00}         
    \end{column}
    \begin{column}{.2\textwidth}
	\includegraphics[width=2cm]{Figs/git/Sourcetree} \newline  \newline  \newline 
    \end{column}
  \end{columns}
\end{frame}


\section{Using git}


\begin{frame}[fragile]{How to use git}
  \begin{columns}[T]
    \begin{column}{.8\textwidth}
  	\textbf{This will assume use of githithub and gitdesktop\\}
	\url{https://desktop.github.com} \newline  \newline 
	Following is a bref dmonstration of the interface
    \end{column}
    \begin{column}{.2\textwidth}
	\includegraphics[width=2cm]{Figs/git/Octocat} \newline  \newline  \newline 
	\includegraphics[width=2cm]{Figs/git/gitdesktop}
    \end{column}
  \end{columns}
\end{frame}


\begin{frame}[fragile]{Adding a new repostry}
  \begin{columns}[T]
    \begin{column}{.12\textwidth}
    \end{column}
    \begin{column}{.88\textwidth}
	\includegraphics[width=10cm]{Figs/GHD/outline_02}
    \end{column}
  \end{columns}
\end{frame}



\begin{frame}[fragile]{Adding a new repostry}
  \begin{columns}[T]
    \begin{column}{.12\textwidth}
    \end{column}
    \begin{column}{.88\textwidth}
	\includegraphics[width=10cm]{Figs/GHD/outline_01}
    \end{column}
  \end{columns}
\end{frame}


\begin{frame}[fragile]{Adding a new repostry}
  \begin{columns}[T]
    \begin{column}{.12\textwidth}
	\small This makes a totally new repository (with no `.git' folder)
    \end{column}
    \begin{column}{.88\textwidth}
	\includegraphics[width=10cm]{Figs/GHD/menu_02}
    \end{column}
  \end{columns}
\end{frame}

\begin{frame}[fragile]{Adding a new repostry}
  \begin{columns}[T]
    \begin{column}{.12\textwidth}
	\small addes an existing repository (with a `.git' folder)
    \end{column}
    \begin{column}{.88\textwidth}
	\includegraphics[width=10cm]{Figs/GHD/menu_03}
    \end{column}
  \end{columns}
\end{frame}

\begin{frame}[fragile]{Adding a new repostry}
  \begin{columns}[T]
    \begin{column}{.12\textwidth}
	\small Copys an online repository to the PC
    \end{column}
    \begin{column}{.88\textwidth}
	\includegraphics[width=10cm]{Figs/GHD/menu_04}
    \end{column}
  \end{columns}
\end{frame}

\begin{frame}[fragile]{Adding a new repostry}
  \begin{columns}[T]
    \begin{column}{.2\textwidth}
	\small new repository menu
    \end{column}
    \begin{column}{.8\textwidth}
	\includegraphics[width=8cm]{Figs/GHD/new_repo_00}
    \end{column}
  \end{columns}
\end{frame}


\begin{frame}[fragile]{Adding a new repostry}
  \begin{columns}[T]
    \begin{column}{.2\textwidth}
    \end{column}
    \begin{column}{.8\textwidth}
	\includegraphics[width=8cm]{Figs/GHD/new_repo_01}
    \end{column}
  \end{columns}
\end{frame}

\begin{frame}[fragile]{Adding a new repostry}
  \begin{columns}[T]
    \begin{column}{.2\textwidth}
    \end{column}
    \begin{column}{.8\textwidth}
	\includegraphics[width=8cm]{Figs/GHD/new_repo_02}
    \end{column}
  \end{columns}
\end{frame}

\begin{frame}[fragile]{Adding a new repostry}
  \begin{columns}[T]
    \begin{column}{.2\textwidth}
	\small The folder \textbf{in which} to store the git reposity folder.
    \end{column}
    \begin{column}{.8\textwidth}
	\includegraphics[width=8cm]{Figs/GHD/new_repo_03}
    \end{column}
  \end{columns}
\end{frame}

\begin{frame}[fragile]{Adding a new repostry}
  \begin{columns}[T]
    \begin{column}{.2\textwidth}
	\small This will inishalise the repositry with a readme containing the discription.
    \end{column}
    \begin{column}{.8\textwidth}
	\includegraphics[width=8cm]{Figs/GHD/new_repo_04}
    \end{column}
  \end{columns}
\end{frame}

\begin{frame}[fragile]{Adding a new repostry}
  \begin{columns}[T]
    \begin{column}{.2\textwidth}
	\small This is a file called `.gitignore' which lists files \& folders to not version control
    \end{column}
    \begin{column}{.8\textwidth}
	\includegraphics[width=8cm]{Figs/GHD/new_repo_05}
    \end{column}
  \end{columns}
\end{frame}


\begin{frame}[fragile]{Using the repostry}
	\includegraphics[width=12cm]{Figs/GHD/outline_02}
\end{frame}


\begin{frame}[fragile]{Using the repostry}
	\small Now we have a inshal commit of the readme done
	\includegraphics[width=12cm]{Figs/GHD/outline_09}
\end{frame}


\begin{frame}[fragile]{Using the repostry}
	\small We can carry on making local commits or publish the repostry to github
	\includegraphics[width=12cm]{Figs/GHD/outline_03}
\end{frame}

\begin{frame}[fragile]{Using the repostry}
	\small We can publish to github
	\includegraphics[width=12cm]{Figs/GHD/init_00}
\end{frame}

\begin{frame}[fragile]{Using the repostry}
	\small Make branches
	\includegraphics[width=12cm]{Figs/GHD/outline_04}
\end{frame}

\begin{frame}[fragile]{Using the repostry}
	\small swich to another reposotory
	\includegraphics[width=12cm]{Figs/GHD/outline_05}
\end{frame}

\begin{frame}[fragile]{Using the repostry}
	\small look at the history of this repository
	\includegraphics[width=12cm]{Figs/GHD/outline_06}
\end{frame}

\begin{frame}[fragile]{Using the repostry}
	\small look at the history of this repository
	\includegraphics[width=12cm]{Figs/GHD/hist_00}
\end{frame}


\begin{frame}[fragile]{Using the repostry}
	\small we can carry on making local commits or or publish the repostry to github
	\includegraphics[width=12cm]{Figs/GHD/outline_07}
\end{frame}

\begin{frame}[fragile]{Using the repostry}
	\small we can carry on making local commits or or publish the repostry to github
	\includegraphics[width=12cm]{Figs/GHD/outline_08}
\end{frame}

{%
\setbeamertemplate{frame footer}{(xkcd)}
{\setbeamercolor{palette primary}{fg=gray, bg=gray}
\begin{frame}[standout]
  \includegraphics[width=6.5cm]{Figs/xkcd}
\end{frame}
}


\appendix

%\begin{frame}[fragile]{Backup slides}
%  Sometimes, it is useful to add slides at the end of your presentation to
%  refer to during audience questions.

%  The best way to do this is to include the \verb|appendixnumberbeamer|
%  package in your preamble and call \verb|\appendix| before your backup slides.

%  \themename will automatically turn off slide numbering and progress bars for
%  slides in the appendix.
%\end{frame}

\begin{frame}[allowframebreaks]{References}

  \bibliography{Mendeley}
  \bibliographystyle{apalike}

\end{frame}

\end{document}
